\documentclass[10pt,a4paper]{article}
\usepackage[T1]{fontenc}
\usepackage[utf8]{inputenc}
\usepackage{lmodern,amsmath,amssymb,booktabs,longtable,array,geometry,microtype,xcolor,url}
\geometry{margin=20mm}
\setlength{\parskip}{3pt}
\setlength{\parindent}{0pt}
\newcommand{\code}[1]{\texttt{#1}}
\newcommand{\rh}{\rho_{AB}}
\title{XXZ Boundary Probes: Exact Settings and Phantom/Helix Preparations}
\author{}
\date{22 September 2026}
\begin{document}
\maketitle
\begin{abstract}
Definitions
\end{abstract}

\section{Hilbert space and Hamiltonians}
The open system has $L$ chain spins and two external probes, ordered $A,1,\ldots,L,B$; $L=8$ gives 10 spins. The probe computational basis is $(|00\rangle,|01\rangle,|10\rangle,|11\rangle)$, with $\sigma^z|0\rangle=+|0\rangle$ and $\sigma^z|1\rangle=-|1\rangle$. Write $\sigma_k^\alpha$ for the Pauli operator in direction $\alpha\in\{x,y,z\}$ acting on spin $k$, and define $\sigma_k^\pm=(\sigma_k^x\pm i\sigma_k^y)/2$. The source sets $\hbar=1$. Its chain and probe Hamiltonians are
\begin{align}
H_C(\Delta)&=\frac{J_{xy}}4\sum_{j=1}^{L-1}\left(
\sigma_j^x\sigma_{j+1}^x+\sigma_j^y\sigma_{j+1}^y
+\Delta\,\sigma_j^z\sigma_{j+1}^z\right)
 +\frac w2\sum_{j=1}^L \sigma_j^z+\frac{e_d}2\sigma_d^z,
& J_z&=J_{xy}\Delta,\\
H_Q&=\frac{\omega_A}{2}\sigma_A^z+\frac{\omega_B}{2}\sigma_B^z,\\
H_{\mathrm{int}}^Z&=\frac{g_A}{4}\sigma_A^z\sigma_1^z
+\frac{g_B}{4}\sigma_L^z\sigma_B^z,\\
H_{\mathrm{int}}^{\mathrm{Exchange}}
&=g_A(\sigma_A^+\sigma_1^-+\sigma_A^-\sigma_1^+)
+g_B(\sigma_L^+\sigma_B^-+\sigma_L^-\sigma_B^+)\\
&=\frac{g_A}{2}(\sigma_A^x\sigma_1^x+\sigma_A^y\sigma_1^y)
+\frac{g_B}{2}(\sigma_L^x\sigma_B^x+\sigma_L^y\sigma_B^y).
\end{align}
The \code{FullHamiltonian} is $H_C+H_Q+H_{\mathrm{int}}^{\mathrm{model}}$, with $H_C$ extended by the probe identities. The code takes the real part of the assembled matrix. It diagonalizes sectors of fixed total $\sigma^z$ and applies unitary phases $e^{-iEt}$; the \code{Z} engine instead evolves four conditional chain Hamiltonians and probe phases. The defect field $e_d\sigma_d^z/2$ is a Hamiltonian term; it is separate from the \emph{initial-state} helix phase defect below.

\section{Exact configuration}
\begin{center}\small
\begin{tabular}{@{}p{34mm}p{50mm}p{73mm}@{}}
\toprule
Source setting & Final \code{Z} value & Final \code{Exchange} value\\\midrule
\code{RunZ}, \code{RunExchange} & \code{True} & \code{True}\\
\code{deltaListZ}, \code{deltaListExchange} & $\{0.01,0.5,1.,1.5,10.\}$ & $\{-0.5,0.,0.5\}$ (later assignment overrides $\{0.01,0.5,1.,1.5,10.\}$)\\
\code{L}, total spins & $8$, $10$ & $8$, $10$\\
\code{Jxy}, $J_z$ & $1.$, $\Delta$ & $1.$, $\Delta$\\
\code{w}, \code{ed}, \code{DefectSite} & $0.$, $0.$, $3$ (field off) & $0.$, $0.$, $3$ (field off)\\
\code{omegaA}, \code{omegaB} & $1.$, $1.$ & $1.$, $1.$\\
\code{gA}, \code{gB} & $0.1$, $0.1$ & $0.1$, $0.1$\\
\code{tMax}, \code{dt} & $1000.$, $0.1$ & $1000.$, $0.1$\\
\code{RunLabel} & \code{run06} & \code{exchange\_helix\_fixed\_q}\\
\code{ChainStateLabel} & \emph{Matched helix with central phase defect} & \emph{Fixed transverse helix, q = Pi/3}\\
\code{ProbeStateLabel} & \emph{|+>\_A tensor |+>\_B} & same\\
\code{GroundTolerance} & $10^{-10}$ & $10^{-10}$\\
\code{ValidationTolerance} & $10^{-8}$ & $10^{-8}$\\
\code{InverseTolerance} & $10^{-8}$ & $10^{-8}$\\
\code{CheckDivisibility} & \code{False} & \code{False}\\
\code{SaveData} & \code{True} & \code{True}\\
\code{PlotChopTolerance} & $10^{-10}$ & $10^{-10}$\\
\code{OutputRoot} & \multicolumn{2}{p{126mm}@{}}{\code{FileNameJoin[\{BaseDirectory[0.], "xxz\_boundary\_results"\}]} (script's directory when loaded from a file)}\\
\bottomrule
\end{tabular}
\end{center}

\section{Initial probe, chain, and total states}
Both active probe functions return $(1,1,1,1)/2=|+\rangle_A\otimes|+\rangle_B$ where $|+\rangle=(|0\rangle+|1\rangle)/\sqrt2$. Thus $\rho_{AB}(0)=|++\rangle\langle++|$, whose 16 matrix entries equal $1/4$ in the stated basis. At each $\Delta$ the chain input is generated separately. In the physical site order $A,1,\ldots,L,B$, the initial state factorizes as
\begin{align}
\rho_{A,C,B}(0)
  &=\rho_A(0)\otimes\rho_C(0)\otimes\rho_B(0),\\
\rho_A(0)&=|+\rangle_A\langle+|,
&\rho_B(0)&=|+\rangle_B\langle+|,\\
\rho_C(0)&=|\Psi_C(\Delta)\rangle\langle\Psi_C(\Delta)|,
&\rho_{AB}(0)&=\rho_A(0)\otimes\rho_B(0).
\end{align}
The full evolution is unitary, and the reported two-probe state is
\begin{equation}
\rho_{AB}(t)=\operatorname{Tr}_C\!\left[
e^{-iHt}\rho_{A,C,B}(0)e^{iHt}
\right].
\end{equation}



\section{Exact names and definitions of saved trajectory quantities}
\enlargethispage{3\baselineskip}
Every sampled time produces one \code{DensityObservables} association. Below $r=\rho_{AB}(t)$, $a=\operatorname{Tr}_B r$, $b=\operatorname{Tr}_A r$, $r_0=\rho_{AB}(0)$, and $\langle O\rangle=\operatorname{Re}\operatorname{Tr}(Or)$ for Hermitian probe operators. The script's entropy \code{VNEntropy} uses positive eigenvalues of the Hermitian part and the natural logarithm.
\begin{center}\small
\begin{longtable}{@{}p{43mm}p{113mm}@{}}
\toprule
Association key & Exact meaning in this density workflow\\\midrule
\endfirsthead
\toprule Association key & Exact meaning in this density workflow\\\midrule\endhead
\code{Delta}, \code{Time} & Run anisotropy $\Delta$ and sampled time $t$.\\
\code{RhoAB}, \code{RhoA}, \code{RhoB} & $r$, $a$, $b$; $4\times4$, $2\times2$, $2\times2$ reduced density matrices.\\
\code{Purity} & $\operatorname{Re}\operatorname{Tr}(r^2)$.\\
\code{EntropyA}, \code{EntropyB}, \code{EntropyAB} & $S(a)$, $S(b)$, $S(r)$ where $S(\rho)=-\sum_{\lambda>0}\lambda\ln\lambda$ (nats).\\
\code{MutualInfo} & $S(a)+S(b)-S(r)$, probe mutual information.\\
\code{Concurrence} & $\max(0,\lambda_1-\lambda_2-\lambda_3-\lambda_4)$: Wootters concurrence; $\lambda_k$ are descending nonnegative square roots from the usual spin-flipped two-qubit construction, with small negative eigenvalues clipped in the implementation.\\
\code{Negativity} & Sum of absolute values of negative eigenvalues of the Hermitian part of $r^{T_B}$; for a normalized state, $(\|r^{T_B}\|_1-1)/2$.\\
\code{InitialOverlap} & $\operatorname{Re}\operatorname{Tr}(r_0r)$; with this pure probe input, $\langle++|r|++\rangle$. It is probe-input survival probability, not a full-chain Loschmidt echo.\\
\code{MagnetizationA}, \code{MagnetizationB} & $\langle \sigma_A^z\rangle$, $\langle \sigma_B^z\rangle$ (Pauli convention, not divided by 2).\\
\code{ConnectedZZ} & $\langle \sigma_A^z\sigma_B^z\rangle-\langle \sigma_A^z\rangle\langle \sigma_B^z\rangle$.\\
\code{ConnectedXX} & $\langle \sigma_A^x\sigma_B^x\rangle-\langle \sigma_A^x\rangle\langle \sigma_B^x\rangle$. Both connected quantities correlate the \emph{probes}.\\
\code{TraceError} & $|\operatorname{Tr}r-1|$.\\
\code{HermiticityError} & $\|r-r^\dagger\|_F$.\\
\code{MinEigenvalue} & Minimum real eigenvalue of $(r+r^\dagger)/2$.\\
\bottomrule
\end{longtable}
\end{center}
For \code{Z}, \code{RhoAB} is the entrywise product of $r_0$ and four-condition chain overlap factors $F_{ij}(t)$; its diagonal probe populations stay fixed under this pure-dephasing model. For \code{Exchange}, \code{RhoAB} is obtained by evolving the full 10-spin probe-chain state and tracing out the chain, allowing population transfer. The six independent pairs $(00,01)$, $(00,10)$, $(00,11)$, $(01,10)$, $(01,11)$, $(10,11)$ are used for the real, imaginary and absolute coherence plots. 




\end{document}
